
\documentclass[../../../physics-standard_answers_all.tex]{subfiles}

\begin{document}

\setlength{\abovedisplayskip}{2mm}
\setlength{\belowdisplayskip}{1mm}

\begin{enumerate}

    \item Yについての電荷保存則より，
        %
        \begin{align*}
            -x_1 + y_1 = 0 \;.
        \end{align*}
        %
        キルヒホッフ則より，
        %  
        \begin{align*}
            \frac{x_1}{C_0} + \frac{y_1}{C_0/2} = 2V_0
            \qquad \text{∴} \; x_1 = y_1
            = \uwave{\frac{2}{3}C_0 V_0} \;.
        \end{align*}
        %
        すると，
        %
        \begin{align*}
            & E_1 = \frac{x_1}{C_0 d} = \uwave{\frac{2V_0}{3d}} \;, \\
            & \phi_1 = \uwave{2V_0} \;,
            \quad \theta_1 = \frac{y_2}{C_0/2} = \uwave{\frac{4}{3}V_0} \;.
        \end{align*}
        
    \item 電荷保存則より，Xの電荷は変わらない．つまり，
        %
        \begin{align*}
            x_2 = x_1 = \uwave{\frac{2}{3}C_0 V_0} \;.
        \end{align*}
        %
        キルヒホッフ則より，
        %
        \begin{align*}
            \frac{y_2}{C_0/2} = V_0
            \qquad \text{∴} \; y_2 = \uwave{\frac{1}{2}C_0 V_0} \;.
        \end{align*}
        %
        すると，
        %
        \begin{align*}
            & E_2 = \frac{x_2}{C_0 d} = \uwave{\frac{2V_0}{3d}} \;, \\
            & \theta_2 = \uwave{V_0} \;,
            \quad \phi_2 = \theta_2 + \frac{x_2}{C_0}
            = \uwave{\frac{5}{3}V_0} \;.
        \end{align*}

    \item Yについての電荷保存則より，
        %
        \begin{align*}
            -x_3 + y_3 = -x_2 + y_2
            \qquad \text{∴} \; -x_3 + y_3 = -\frac{1}{6}C_0 V_0 \;.
        \end{align*}
        %
        キルヒホッフ則より，
        %  
        \begin{align*}
            & \frac{x_3}{C_0} + \frac{y_3}{C_0/2} = 2V_0 \\
            & \text{∴} \; x_3 = \uwave{\frac{7}{9}C_0 V_0} \;,
            \quad y_3 = \uwave{\frac{11}{18}C_0 V_0} \;.
        \end{align*}
        %
        すると，
        %
        \begin{align*}
            & E_3 = \frac{x_3}{C_0 d} = \uwave{\frac{7V_0}{9d}} \;, \\
            & \phi_3 = \uwave{2V_0} \;,
            \quad \theta_3 = \frac{y_3}{C_0/2} = \uwave{\frac{11}{9}V_0} \;.
        \end{align*}

    \item X下面の電荷を$x_4$，Y下面の電荷を$y_4$とすると，Yについての電荷保存則より，
        %
        \begin{align*}
            -x_4 + y_4 = -x_3 + y_3
            \qquad \text{∴} \; -x_4 + y_4 = -\frac{1}{6}C_0 V_0 \;.
        \end{align*}
        %
        キルヒホッフ則より，
        %  
        \begin{align*}
            & \frac{x_4}{C_0/2} + \frac{y_4}{C_0} = 2V_0 \\
            & \text{∴} \; x_4 = \frac{13}{18}C_0 V_0 \;,
            \quad y_4 = \frac{5}{9}C_0 V_0 \;.
        \end{align*}
        %
        $\Delta Q$は，Xの電荷の変化量に等しいから，
        %
        \begin{align*}
            \Delta Q = x_4 - x_3 = \uwave{-\frac{1}{18}C_0 V_0} \;.
        \end{align*}
        
\end{enumerate}

\end{document}
